\section{Weekly Summary}

\begin{tcolorbox}[colback=yellow!10,colframe=black!50,title=AI-Generated Content]
\noindent The summary below was written by Claude (Anthropic) from that week's regression outputs and series descriptions. It has not been reviewed or fact-checked by a human, and should be read as an automated, informal recap -- not analysis.
\end{tcolorbox}

Welcome back to the only weekly column brave enough to pair economic time series like a blindfolded matchmaker and then act surprised when they don't get along. As always, the disclaimer that keeps me out of statistical jail: these are randomly chosen series with zero reason to be related, and any correlation you see here is the universe's idea of a practical joke. Onward.

Monday kicked things off with New England commercial electricity prices versus a price index for medical expenditures on, and I quote, "symptoms, signs, and ill-defined conditions." The raw regression came out screaming, R-squared of 0.93, a p-value with more zeros than my bank account. Absolutely thrilling! Right up until you detrend it, at which point the whole thing evaporates into an R-squared of 0.02 and a p-value of basically "who knows." Classic trend mirage. Two things went up over time, held hands, and pretended it meant something. It did not. Tuesday gave us a rarer treat: Montana housing listings in Flathead County against U.S. imports of maintenance and repair services. Weaker to begin with (raw R-squared of 0.12), but here's the twist -- the relationship actually got STRONGER after detrending, up to 0.19 and still very significant. I have no idea why active real estate listings in one Montana county would relate to national repair-service imports, and I refuse to invent a theory, but I'll grudgingly tip my hat: it survived the detrending guillotine, which is more than most contestants manage.

Wednesday was the week's overachiever. Wholesale-trade payrolls versus the Dover, Delaware house price index posted a raw R-squared of 0.73, and -- brace yourself -- it held onto 0.73 after detrending too, p-value still buried in the scientific-notation basement. Two independent series that both drift AND wiggle together? Suspicious. Probably still coincidence dressed in a nice suit, because payrolls and house prices both quietly track "the economy doing economy things," but as far as this project goes, that's a genuine standout worth pointing at. Thursday, meanwhile, gave us Southwest central-bank GDP against a gloriously named discontinued bank-assets series, where the raw result was borderline (p of 0.05, a coward's significance) but the detrended version firmed up to a real negative relationship. Interpret at your own peril.

Friday paired education-and-communication consumer prices with total tax exemptions under age 65 in New Mexico, and this one actually doubled down: R-squared jumped from 0.32 raw to 0.51 detrended. Two more series conspiring behind the trend's back. Then the weekend arrived to humble us all. Saturday's "shares of gross domestic income" versus a Peabody, Massachusetts house price index produced a detrended coefficient of roughly 57 trillion with a p-value of 0.82 and warnings about a singular matrix -- which is math's polite way of saying "please stop." And Sunday's employment cost index versus U.S. JPY intervention transactions returned nothing but zeros and NaNs, the regression equivalent of the machine unplugging itself and walking out. Fitting end to the week: a strong reminder that spurious is the house specialty, the occasional survivor is just noise with good posture, and I'll be right back here next week to do it all again.
